\section*{Author's foreword}
{\small\sffamily\color{copper}FIRST-PERSON DRAFT FOR ARYAN AYYAR'S REVIEW}\par\medskip
My interest in Kai Tak began with a childhood fascination with aircraft. A large passenger aeroplane turning above Hong Kong seemed to bring two very different scales into the same scene: the machine in flight and the streets below it. I wanted to understand how they fitted together.

Richard Hammond's \emph{Engineering Connections} encouraged that curiosity. Understanding a machine's workings made it more absorbing to watch. The imagery and music of \emph{Enter the Dragon} were another part of my early interest in Hong Kong. Later, studying at Hong Kong Polytechnic University gave the city a place in my own life. My time there was after Kai Tak had closed; the airport I arrived at was Chek Lap Kok.

Hong Kong matters to me personally, beyond this research question. The question itself is strong enough: how could an aircraft of this size turn towards that runway and still complete its landing?

I approach it as a financial economist and mathematician. The tools are familiar; the engineering invites careful learning. Following the turn leads to forces, speed and descent. Following the aircraft to a full stop brings the runway into the calculation. Returning to old photographs brings the city back into view. That movement between a place, a physical question and a mathematical explanation is what drew me to the project, and what I hope will draw a reader through it.
